% F085｜逐类防错；明确“模型”不是唯一解析式。
\FReserve{3.5cm}
\begin{trapbox}
\textbf{\textcolor{CTrap}{易错 1｜把识别线索当成唯一解析式}}

$f(x+y)=f(x)f(y)$ \textbf{不自动等于} $f(x)=a^x$：它还允许 $f\equiv0$。$f(x+y)=f(x)+f(y)$ \textbf{不自动等于} $kx$；若要唯一写出常见表达式，还需连续性等附加条件。

\textbf{\textcolor{CTrap}{易错 2｜忘了模型有自己的定义域}}

对数型、幂函数型先在 $x,y>0$ 讨论；正切、余切不能在自身无定义的位置使用加法公式。不能把某个式子当作对全体实数恒成立。
\end{trapbox}

\FReserve{3.5cm}
\begin{trapbox}
\textbf{\textcolor{CTrap}{易错 3｜指数型漏掉零函数}}

由 $f(x)=f(x)f(0)$，只有\textbf{存在某个 $x$ 使 $f(x)\ne0$} 时，才能推出 $f(0)=1$。若 $f\equiv0$，则 $f(0)=0$。
% F085｜真实数学图；代表函数不构成抽象解的唯一性证明。
\begin{center}
\begin{tikzpicture}
\begin{axis}[width=12.4cm,height=5.1cm,axis lines=middle,
 xmin=-2.5,xmax=2.6,ymin=-0.5,ymax=4.8,
 xtick={-2,-1,0,1,2},xticklabels=\empty,ytick={0,1,2,4},xlabel={$t$},ylabel={$f(t)$},
 title={指数型的两个分支：非零代表与零函数},title style={font=\small,yshift=3pt},
 tick label style={font=\small},label style={font=\small},
 ylabel style={at={(axis description cs:0,1)},anchor=south east,xshift=-18pt},
 clip=false,grid=major,grid style={black!8},enlargelimits=false,
 every axis plot/.append style={thick},samples=161]
\addplot[CExplanation,very thick,domain=-2:2.1] {2^x};
\addplot[orange!80!black,dashed,domain=-2:2.1] {0};\addplot[only marks,mark=*,mark size=2pt,orange!80!black] coordinates {(0,1) (0,0)};

\node[font=\small,anchor=north,yshift=-3pt] at (axis cs:-2,-0.5) {$-2$};
\node[font=\small,anchor=north,yshift=-3pt] at (axis cs:-1,-0.5) {$-1$};
\node[font=\small,anchor=north,yshift=-3pt] at (axis cs:1,-0.5) {$1$};
\node[font=\small,anchor=north,yshift=-3pt] at (axis cs:2,-0.5) {$2$};

\end{axis}
\end{tikzpicture}
\par\small 蓝实线 $2^t$ 与橙虚线 $0$ 都满足 $f(x+y)=f(x)f(y)$，在 $0$ 处的值分别为 $1,0$；先排除零函数才能约去非零值。
\end{center}


\textbf{\textcolor{CTrap}{易错 4｜对数型和幂函数型混淆}}
\[
\text{对数： }f(x^n)=nf(x),\qquad
\text{幂： }f(x^n)=[f(x)]^n.
\]
如果题目只有 $f(2)$，求 $f(8)$ 时要区分 $3f(2)$ 与 $f(2)^3$。
% F085｜真实数学图；代表函数不构成抽象解的唯一性证明。
\begin{center}
\begin{minipage}{.49\linewidth}\centering
\begin{tikzpicture}
\begin{axis}[width=6.1cm,height=4.8cm,axis lines=middle,
 xmin=-0.2,xmax=8.9,ymin=-4,ymax=30,
 xtick={0,2,4,8},xticklabels=\empty,ytick={0,3,6,9,27},xlabel={$t$},ylabel={$f(t)$},
 title={同尺度对照：$3\log_2t$},title style={font=\small,yshift=3pt},
 tick label style={font=\small},label style={font=\small},
 ylabel style={at={(axis description cs:0,1)},anchor=south east,xshift=-18pt},
 clip=false,grid=major,grid style={black!8},enlargelimits=false,
 every axis plot/.append style={thick},samples=161]
\addplot[CExplanation,very thick,domain=.5:8.1] {3*ln(x)/ln(2)};
\draw[dashed,black!35] (axis cs:2,0) |- (axis cs:0,3);
\draw[dashed,black!35] (axis cs:4,0) |- (axis cs:0,6);
\draw[dashed,black!35] (axis cs:8,0) |- (axis cs:0,9);
\addplot[only marks,mark=*,mark size=2pt,orange!80!black] coordinates {(2,3) (4,6) (8,9)};

\node[font=\small,anchor=north,yshift=-3pt] at (axis cs:2,-4) {$2$};
\node[font=\small,anchor=north,yshift=-3pt] at (axis cs:4,-4) {$4$};
\node[font=\small,anchor=north,yshift=-3pt] at (axis cs:8,-4) {$8$};

\end{axis}
\end{tikzpicture}
\end{minipage}\hfill
\begin{minipage}{.49\linewidth}\centering
\begin{tikzpicture}
\begin{axis}[width=6.1cm,height=4.8cm,axis lines=middle,
 xmin=-0.2,xmax=8.9,ymin=-4,ymax=30,
 xtick={0,2,4,8},xticklabels=\empty,ytick={0,3,6,9,27},xlabel={$t$},ylabel={$f(t)$},
 title={同尺度对照：$t^{\log_2 3}$},title style={font=\small,yshift=3pt},
 tick label style={font=\small},label style={font=\small},
 ylabel style={at={(axis description cs:0,1)},anchor=south east,xshift=-18pt},
 clip=false,grid=major,grid style={black!8},enlargelimits=false,
 every axis plot/.append style={thick},samples=161]
\addplot[CExplanation,very thick,domain=.5:8.1] {x^(ln(3)/ln(2))};
\draw[dashed,black!35] (axis cs:2,0) |- (axis cs:0,3);
\draw[dashed,black!35] (axis cs:4,0) |- (axis cs:0,9);
\draw[dashed,black!35] (axis cs:8,0) |- (axis cs:0,27);
\addplot[only marks,mark=*,mark size=2pt,orange!80!black] coordinates {(2,3) (4,9) (8,27)};

\node[font=\small,anchor=north,yshift=-3pt] at (axis cs:2,-4) {$2$};
\node[font=\small,anchor=north,yshift=-3pt] at (axis cs:4,-4) {$4$};
\node[font=\small,anchor=north,yshift=-3pt] at (axis cs:8,-4) {$8$};

\end{axis}
\end{tikzpicture}
\end{minipage}
\par\small 都过 $(2,3)$，在 $4$ 处却读到 $6$ 和 $9$，在 $8$ 处读到 $9$ 和 $27$。一个点的值不能代替运算结构。
\end{center}

\end{trapbox}

\FReserve{3.5cm}
\begin{trapbox}
\textbf{\textcolor{CTrap}{易错 5｜余弦型不一定就是 $\cos x$}}

$f(x)\equiv1$ 同样满足余弦型方程；
\[
 f(x)=\frac{2^x+2^{-x}}2
\]
也满足它。因此不能没有条件就说函数“必有周期 $2\pi$”或“函数值总在 $[-1,1]$”。

\textbf{\textcolor{CTrap}{易错 6｜正弦型不等于只有正弦函数}}

$f(x)=x$ 也满足正弦型平方差关系；更一般地 $f(x)=cx$（$c$ 为常数）也满足。\textbf{不要把正弦函数的周期直接推广给所有正弦型解。}
% F085｜真实数学图；代表函数不构成抽象解的唯一性证明。
\begin{center}
\begin{minipage}{.49\linewidth}\centering
\begin{tikzpicture}
\begin{axis}[width=6.1cm,height=4.8cm,axis lines=middle,
 xmin=-2.9,xmax=2.9,ymin=-1.5,ymax=3.4,
 xtick={-2,0,2},xticklabels=\empty,ytick={-1,0,1,2,3},xlabel={$t$},ylabel={$f(t)$},
 title={同属余弦型，形状却不同},title style={font=\small,yshift=3pt},
 tick label style={font=\small},label style={font=\small},
 ylabel style={at={(axis description cs:0,1)},anchor=south east,xshift=-18pt},
 clip=false,grid=major,grid style={black!8},enlargelimits=false,
 every axis plot/.append style={thick},samples=161]
\addplot[CExplanation,very thick,domain=-2.5:2.5] {cos(deg(x))};
\addplot[orange!80!black,dashed,domain=-2.5:2.5] {(2^x+2^(-x))/2};\addplot[green!45!black,dotted,domain=-2.5:2.5] {1};
\node[font=\small,anchor=north,yshift=-3pt] at (axis cs:-2,-1.5) {$-2$};
\node[font=\small,anchor=north,yshift=-3pt] at (axis cs:2,-1.5) {$2$};

\end{axis}
\end{tikzpicture}
\end{minipage}\hfill
\begin{minipage}{.49\linewidth}\centering
\begin{tikzpicture}
\begin{axis}[width=6.1cm,height=4.8cm,axis lines=middle,
 xmin=-2.9,xmax=2.9,ymin=-3.1,ymax=3.1,
 xtick={-2,0,2},xticklabels=\empty,ytick={-2,0,2},xlabel={$t$},ylabel={$f(t)$},
 title={同属正弦型：曲线与直线},title style={font=\small,yshift=3pt},
 tick label style={font=\small},label style={font=\small},
 ylabel style={at={(axis description cs:0,1)},anchor=south east,xshift=-18pt},
 clip=false,grid=major,grid style={black!8},enlargelimits=false,
 every axis plot/.append style={thick},samples=161]
\addplot[CExplanation,very thick,domain=-2.5:2.5] {sin(deg(x))};
\addplot[orange!80!black,dashed,domain=-2.5:2.5] {x};
\node[font=\small,anchor=north,yshift=-3pt] at (axis cs:-2,-3.1) {$-2$};
\node[font=\small,anchor=north,yshift=-3pt] at (axis cs:2,-3.1) {$2$};

\end{axis}
\end{tikzpicture}
\end{minipage}
\par\small 左图：蓝实线 $\cos t$，橙虚线 $(2^t+2^{-t})/2$，绿点线 $1$；右图：蓝实线 $\sin t$，橙虚线 $t$。同一模型容纳不同曲线，不能直接断言周期、值域或唯一解析式。
\end{center}

\end{trapbox}

\FReserve{3.5cm}
\begin{trapbox}
\textbf{\textcolor{CTrap}{易错 7｜正切型分母为零还照套}}

只有 $1-f(x)f(y)\ne0$，才能使用
\[
f(x+y)=\frac{f(x)+f(y)}{1-f(x)f(y)}.
\]
对具体 $\tan$，若 $\cos x=0$、$\cos y=0$ 或 $\cos(x+y)=0$，相关函数值可能无定义；\textbf{仅写代数式，不检查输入是否合法，会出错。}

\textbf{\textcolor{CTrap}{易错 8｜余切型随便代入 $0$}}

具体余切函数 $\cot0$ 不存在；且加法型关系的分母 $f(x)+f(y)$ 还不能为零。\textbf{余切型不能照搬指数型“令 $y=0$”的套路。}
\end{trapbox}

\FReserve{3.5cm}
\begin{trapbox}
\textbf{\textcolor{CTrap}{易错 9｜正切和余切分式记反}}
\[
\begin{aligned}
\text{正切：}&\quad\frac{f(x)+f(y)}{1-f(x)f(y)},\\[4pt]
\text{余切：}&\quad\frac{f(x)f(y)-1}{f(x)+f(y)}.
\end{aligned}
\]
\textbf{正切“和在上”，余切“积减一在上”。}

\textbf{\textcolor{CTrap}{易错 10｜正弦型左右两个输入不会反凑}}

$f(m)f(n)$ 要凑 $f(x+y)f(x-y)$，应令
\[
 x=\frac{m+n}2,\quad y=\frac{m-n}2.
\]
不要把“中点”误写成 $m+n$，也不要把“半差”误写成 $m-n$。
\end{trapbox}

\FReserve{3.5cm}
\begin{trapbox}
\textbf{\textcolor{CTrap}{易错 11｜以为所有模型都可推出奇偶或周期}}

正比例型在全体实数上满足加法关系时可以推出奇函数；余弦型不恒为零的解可推出偶函数；正弦型在全体实数上可推出奇函数。\textbf{这些都不能无条件变成某一个固定周期。}

\textbf{\textcolor{CTrap}{易错 12｜只用一个数字验证就断言函数恒等式}}

取 $x=y=1$ 只能帮助发现关系；正式证明某函数满足模型，要按公式对所有允许的输入推导。

\textbf{统一防错顺序：}\textbf{看输入是否允许 $\longrightarrow$ 看分母是否为零 $\longrightarrow$ 认准运算结构 $\longrightarrow$ 再代入特殊值。}
\end{trapbox}

